\section{Results}
\label{sec:results}

\subsection{Sequencing yield and quality control}

The six libraries returned between 29.6 and 41.3 million read pairs, a spread
narrow enough that no library needed downsampling before normalisation.
Table~\ref{tab:qc} gives the per-library summary. Two figures of merit deserve
comment. Mapping rate is uniformly high and varies by less than two percentage
points across libraries, which argues against contamination in any single
flask. Optical duplication, in contrast, is not uniform: NL-2 and NL-3 sit at
19.4 and 21.7 percent against a mean of 11.8 percent for the other four.

\begin{table}[t]
  \centering
  \caption{Sequencing and alignment summary. NR denotes a nitrogen-replete
  replicate and NL a nitrogen-limited one. Duplication is the optical and PCR
  duplicate fraction reported after alignment. Insert size is the median of the
  properly paired distribution.}
  \label{tab:qc}
  \begin{tabular}{lrrrrr}
    \toprule
    Library & Read pairs & Bases $\geq$ Q30 & Duplicates & Mapped & Median insert \\
            & ($10^{6}$) & (\%) & (\%) & (\%) & (bp) \\
    \midrule
    NR-1 & 38.4 & 94.1 & 11.2 & 93.1 & 241 \\
    NR-2 & 41.3 & 93.8 & 12.6 & 92.8 & 236 \\
    NR-3 & 33.1 & 94.4 & 10.9 & 93.4 & 244 \\
    NL-1 & 34.2 & 93.9 & 12.5 & 92.2 & 233 \\
    NL-2 & 31.7 & 92.6 & 19.4 & 91.6 & 228 \\
    NL-3 & 29.6 & 92.1 & 21.7 & 91.4 & 226 \\
    \midrule
    Median & 33.7 & 93.9 & 12.6 & 92.5 & 235 \\
    \bottomrule
  \end{tabular}
\end{table}

Per-cycle base quality (Figure~\ref{fig:quality}) shows the usual decline over
the read, steeper in the second read than the first, but the mean stays above
Q30 for the whole of read one and until cycle 128 of read two. Trimming removed
a median of 4.1 percent of bases. The insert size distribution
(Figure~\ref{fig:inserts}) is unimodal with a median of 235 bp and no adapter
dimer peak below 120 bp, so the elevated duplication in NL-2 and NL-3 is not
a short-fragment artefact.

We retained NL-2 and NL-3 for three reasons. Their duplication is concentrated
in the highest-abundance decile of transcripts, where it barely perturbs the
rank order. Removing duplicates before counting changed the fitted regime
coefficient by less than $0.08$ in $\logfc$ for every transcript that reached
significance. And a leave-one-out reanalysis dropping either library recovered
94.1 and 95.6 percent of the significant set, respectively.

\begin{figure}[t]
  \centering
  \begin{tikzpicture}
    \begin{axis}[
        width=0.82\linewidth, height=6.0cm,
        xlabel={sequencing cycle},
        ylabel={mean Phred quality},
        xmin=0, xmax=152, ymin=24, ymax=40,
        grid=both,
        major grid style={draw=gridline,line width=0.4pt},
        minor grid style={draw=gridline!55,line width=0.3pt},
        tick label style={font=\footnotesize},
        label style={font=\small},
        legend style={font=\footnotesize,at={(0.03,0.04)},anchor=south west,
          draw=gridline,fill=white},
        legend cell align=left,
      ]
      \addplot[draw=none,forget plot,name path=qfloor]
        coordinates {(0,30) (152,30)};
      \addplot[draw=none,forget plot,name path=qceil]
        coordinates {(0,40) (152,40)};
      \addplot[drylab!12,forget plot] fill between[of=qfloor and qceil];

      \addplot[wetlab,line width=1pt,mark=*,mark size=1pt] coordinates {
        (1,32.1) (5,35.4) (10,36.2) (20,36.4) (30,36.3) (40,36.2)
        (50,36.0) (60,35.9) (70,35.7) (80,35.4) (90,35.1) (100,34.7)
        (110,34.2) (120,33.6) (130,32.8) (140,31.9) (150,30.6)
      };
      \addlegendentry{read 1}
      \addplot[accent,line width=1pt,dashed,mark=square*,mark size=1pt]
        coordinates {
        (1,31.2) (5,34.6) (10,35.5) (20,35.6) (30,35.4) (40,35.2)
        (50,34.9) (60,34.6) (70,34.3) (80,33.9) (90,33.4) (100,32.8)
        (110,32.1) (120,31.2) (130,30.2) (140,29.0) (150,27.6)
      };
      \addlegendentry{read 2}

      \draw[draw=muted,line width=0.5pt,dash pattern=on 2pt off 2pt]
        (axis cs:0,30) -- (axis cs:152,30);
      \node[anchor=north west,font=\footnotesize,text=muted]
        at (axis cs:3,29.8) {Q30 acceptance floor};
      \draw[-{Stealth[length=2mm]},draw=accent,line width=0.6pt]
        (axis cs:118,27.6) -- (axis cs:128,29.7);
      \node[anchor=north,font=\footnotesize,text=accent,align=center]
        at (axis cs:112,27.4) {read 2 crosses\\ Q30 at cycle 128};
    \end{axis}
  \end{tikzpicture}
  \caption{Mean base quality by sequencing cycle, pooled over all six
  libraries. The shaded band is the region at or above Q30.}
  \label{fig:quality}
\end{figure}

\begin{figure}[t]
  \centering
  \begin{tikzpicture}
    \begin{axis}[
        width=0.82\linewidth, height=5.6cm,
        ybar,
        bar width=6.5pt,
        xlabel={insert size (bp)},
        ylabel={properly paired fragments (\%)},
        xmin=95, xmax=505, ymin=0, ymax=21,
        ytick={0,5,10,15,20},
        grid=major,
        major grid style={draw=gridline,line width=0.4pt},
        tick label style={font=\footnotesize},
        label style={font=\small},
      ]
      \addplot[draw=wetlab,fill=wetlab!35,line width=0.5pt] coordinates {
        (112,0.6) (137,2.1) (162,6.4) (187,12.8) (212,17.9) (237,18.6)
        (262,14.7) (287,10.1) (312,6.5) (337,4.0) (362,2.4) (387,1.4)
        (412,0.9) (437,0.5) (462,0.3) (487,0.2)
      };
      \draw[draw=accent,line width=0.8pt,dash pattern=on 3pt off 2pt]
        (axis cs:235,0) -- (axis cs:235,19.6);
      \node[anchor=west,font=\footnotesize,text=accent]
        at (axis cs:245,19.2) {median 235 bp};
      \node[anchor=north west,font=\footnotesize,text=muted,align=left]
        at (axis cs:300,15.5)
        {no adapter dimer peak\\ below 120 bp};
    \end{axis}
  \end{tikzpicture}
  \caption{Pooled insert size distribution in 25 bp bins. The fragmentation
  time was chosen to place the mode near 230 bp, which the observed
  distribution matches.}
  \label{fig:inserts}
\end{figure}

\subsection{Differential expression}

Of 11{,}776 transcripts passing the independent filter, 1{,}486 were
differentially expressed at an adjusted $p$ below 0.05 with an absolute
$\logfc$ above 1.0. Induction outweighed repression, 863 to 623.
Figure~\ref{fig:volcano} shows the full result and labels the five transcripts
that carry the largest effects. Table~\ref{tab:de} lists the ten strongest in
each direction.

\begin{figure}[t]
  \centering
  \begin{tikzpicture}
    \begin{axis}[
        width=0.86\linewidth, height=7.4cm,
        xlabel={$\log_{2}$ fold change, limited against replete},
        ylabel={$-\log_{10}$ adjusted $p$},
        xmin=-6.6, xmax=6.6, ymin=-0.6, ymax=31,
        grid=both,
        major grid style={draw=gridline,line width=0.4pt},
        minor grid style={draw=gridline!55,line width=0.3pt},
        tick label style={font=\footnotesize},
        label style={font=\small},
        legend style={font=\footnotesize,at={(0.02,0.03)},anchor=south west,
          draw=gridline,fill=white},
        legend cell align=left,
      ]
      \addplot[only marks,mark=*,mark size=1pt,
               mark options={fill=muted!50,draw=muted!60}] coordinates {
      (-1.73,1.00) (-1.72,0.29) (-1.63,1.07) (-1.58,1.14)
      (-1.58,0.68) (-1.55,0.15) (-1.50,0.63) (-1.45,0.19)
      (-1.40,0.30) (-1.37,0.84) (-1.36,0.13) (-1.32,0.65)
      (-1.14,0.40) (-1.11,0.44) (-1.05,0.39) (-1.04,0.40)
      (-0.97,0.32) (-0.96,0.59) (-0.93,0.74) (-0.78,0.53)
      (-0.75,0.19) (-0.49,1.12) (-0.33,0.79) (-0.31,0.97)
      (-0.29,0.29) (-0.23,1.01) (-0.19,0.28) (-0.08,0.66)
      (-0.04,0.24) (-0.02,0.07) (0.11,0.49) (0.42,0.32)
      (0.50,0.26) (0.61,0.54) (0.61,0.18) (0.63,0.49)
      (0.66,0.24) (0.72,0.73) (0.74,0.29) (0.85,1.16)
      (0.90,0.15) (0.93,0.73) (0.96,0.42) (0.97,0.63)
      (1.02,0.16) (1.24,1.17) (1.25,0.52) (1.26,0.56)
      (1.28,0.63) (1.35,0.64) (1.58,0.32) (1.59,0.32)
      (1.66,0.82) (1.66,0.11) (1.69,0.48) (1.69,0.98)
      (1.72,0.14) (1.82,0.47)
      };
      \addlegendentry{not significant}

      \addplot[only marks,mark=*,mark size=1.3pt,
               mark options={fill=wetlab,draw=wetlab}] coordinates {
      (-3.80,16.08) (-3.74,14.27) (-3.71,9.07) (-3.32,11.90)
      (-3.12,9.06) (-3.12,9.02) (-2.91,13.03) (-2.73,9.28)
      (-2.63,10.73) (-2.56,12.14) (-2.51,11.46) (-2.47,11.45)
      (-2.45,6.60) (-2.45,10.32) (-2.26,6.88) (-2.19,6.10)
      (-2.06,4.73) (-2.00,4.72) (-1.99,5.88) (-1.83,6.94)
      (-1.76,6.28) (-1.69,7.32) (-1.58,5.77) (-1.53,7.17)
      (-1.40,6.57) (-1.33,5.88) (-4.62,21.40) (-5.31,26.80)
      };
      \addlegendentry{repressed, 623 transcripts}

      \addplot[only marks,mark=*,mark size=1.3pt,
               mark options={fill=accent,draw=accent}] coordinates {
      (1.21,4.44) (1.29,3.44) (1.32,5.61) (1.41,3.55)
      (1.43,4.22) (1.44,5.14) (1.48,7.37) (1.76,9.11)
      (1.81,4.92) (1.89,8.28) (1.98,7.89) (1.99,4.52)
      (2.05,8.19) (2.06,7.81) (2.09,7.58) (2.22,10.55)
      (2.24,10.91) (2.34,7.84) (2.39,6.72) (2.40,7.52)
      (2.41,11.09) (2.56,10.91) (2.59,11.27) (2.80,11.11)
      (2.82,11.71) (2.89,5.98) (3.00,6.58) (3.24,6.80)
      (3.43,15.47) (3.51,17.95) (4.01,12.52) (4.04,20.03)
      (4.83,24.60) (5.42,28.90) (5.71,27.10)
      };
      \addlegendentry{induced, 863 transcripts}

      \draw[draw=muted,line width=0.5pt,dash pattern=on 2pt off 2pt]
        (axis cs:-6.6,1.301) -- (axis cs:6.6,1.301);
      \draw[draw=muted,line width=0.5pt,dash pattern=on 2pt off 2pt]
        (axis cs:-1,-0.6) -- (axis cs:-1,31);
      \draw[draw=muted,line width=0.5pt,dash pattern=on 2pt off 2pt]
        (axis cs:1,-0.6) -- (axis cs:1,31);

      \node[anchor=east,font=\footnotesize,text=accent]
        at (axis cs:5.30,28.90) {NRT2.1};
      \node[anchor=east,font=\footnotesize,text=accent]
        at (axis cs:5.60,27.10) {AMT1.3};
      \node[anchor=east,font=\footnotesize,text=accent]
        at (axis cs:4.72,24.60) {GSIII};
      \node[anchor=west,font=\footnotesize,text=wetlab]
        at (axis cs:-5.20,26.80) {LHCF7};
      \node[anchor=west,font=\footnotesize,text=wetlab]
        at (axis cs:-4.51,21.40) {PSBS-like};
    \end{axis}
  \end{tikzpicture}
  \caption{Volcano plot of the regime contrast. Dashed lines mark the
  significance threshold at an adjusted $p$ of 0.05 and the fold change
  threshold at $\lvert \logfc \rvert = 1$. The five labelled transcripts are
  discussed in Section~\ref{sec:discussion}.}
  \label{fig:volcano}
\end{figure}

\begin{table}[t]
  \centering
  \small
  \caption{Ten most induced and ten most repressed transcripts. Base mean is
  the size-factor normalised mean count across all six libraries. Fold changes
  are shrunken estimates and adjusted $p$ values follow the Benjamini and
  Hochberg procedure.}
  \label{tab:de}
  \begin{tabular}{llrrr}
    \toprule
    Locus & Annotated product & Base mean & $\logfc$ & Adjusted $p$ \\
    \midrule
    \multicolumn{5}{l}{\emph{Induced under nitrogen limitation}} \\
    TP-04122 & Nitrate transporter NRT2.1        & 4\,812 & $+5.42$ & $1.3 \times 10^{-29}$ \\
    TP-11907 & Ammonium transporter AMT1.3       & 2\,146 & $+5.71$ & $7.9 \times 10^{-28}$ \\
    TP-02318 & Glutamine synthetase III          & 9\,033 & $+4.83$ & $2.5 \times 10^{-25}$ \\
    TP-07741 & Formamidase                       & 1\,204 & $+4.04$ & $9.3 \times 10^{-21}$ \\
    TP-00983 & Nitrate reductase                 & 6\,517 & $+3.51$ & $1.1 \times 10^{-18}$ \\
    TP-13650 & Urea transporter DUR3             &   784 & $+3.43$ & $3.4 \times 10^{-16}$ \\
    TP-05402 & Carbamoyl phosphate synthase      & 1\,873 & $+4.01$ & $3.0 \times 10^{-13}$ \\
    TP-09115 & Purine catabolism allantoinase    &   612 & $+2.82$ & $1.9 \times 10^{-12}$ \\
    TP-06238 & Amino acid permease               & 1\,441 & $+2.59$ & $5.4 \times 10^{-12}$ \\
    TP-01760 & Aspartate aminotransferase        & 3\,290 & $+2.41$ & $8.1 \times 10^{-12}$ \\
    \midrule
    \multicolumn{5}{l}{\emph{Repressed under nitrogen limitation}} \\
    TP-03017 & Fucoxanthin protein LHCF7         & 18\,405 & $-5.31$ & $1.6 \times 10^{-27}$ \\
    TP-03022 & Photosystem II subunit S-like     & 12\,760 & $-4.62$ & $4.0 \times 10^{-22}$ \\
    TP-08844 & Fucoxanthin protein LHCF3         &  9\,918 & $-3.80$ & $8.3 \times 10^{-17}$ \\
    TP-10236 & Ribosomal protein L4              &  7\,602 & $-3.74$ & $5.4 \times 10^{-15}$ \\
    TP-10241 & Ribosomal protein S8              &  6\,845 & $-3.71$ & $8.5 \times 10^{-10}$ \\
    TP-04490 & Ferredoxin                        &  5\,331 & $-3.32$ & $1.3 \times 10^{-12}$ \\
    TP-12008 & Phycobilin lyase-like             &  2\,207 & $-3.12$ & $8.7 \times 10^{-10}$ \\
    TP-00512 & Plastid ATP synthase subunit b    &  8\,116 & $-2.91$ & $9.3 \times 10^{-14}$ \\
    TP-07019 & Chlorophyll synthase              &  3\,458 & $-2.73$ & $5.2 \times 10^{-10}$ \\
    TP-02904 & Ribosomal protein L11             &  6\,031 & $-2.63$ & $1.9 \times 10^{-11}$ \\
    \bottomrule
  \end{tabular}
\end{table}

\subsection{Functional grouping}

Table~\ref{tab:enrichment} collapses the significant set into the annotation
categories that survive a hypergeometric test with the same false discovery
threshold. Nitrogen acquisition and assimilation are the dominant induced
groups. The repressed side is narrower and is carried almost entirely by
light harvesting and cytosolic translation.

\begin{table}[t]
  \centering
  \small
  \caption{Enriched functional categories. Counts are transcripts in the
  significant set against the number annotated to the category in the
  reference. The enrichment ratio is the observed proportion divided by the
  background proportion.}
  \label{tab:enrichment}
  \begin{tabular}{llrrrr}
    \toprule
    Direction & Category & In set & Annotated & Ratio & Adjusted $p$ \\
    \midrule
    Induced   & Nitrate and nitrite transport   &  19 &  24 & 6.2 & $2.1 \times 10^{-14}$ \\
    Induced   & Glutamine and glutamate cycle   &  27 &  41 & 5.2 & $6.6 \times 10^{-13}$ \\
    Induced   & Urea and purine catabolism      &  22 &  38 & 4.6 & $4.0 \times 10^{-9}$ \\
    Induced   & Vacuolar proteolysis            &  34 &  81 & 3.3 & $1.2 \times 10^{-6}$ \\
    Induced   & Lipid droplet biogenesis        &  18 &  47 & 3.0 & $8.4 \times 10^{-5}$ \\
    Repressed & Light harvesting complex        &  31 &  44 & 5.6 & $3.3 \times 10^{-17}$ \\
    Repressed & Cytosolic ribosome              &  63 & 118 & 4.2 & $9.1 \times 10^{-16}$ \\
    Repressed & Photosynthetic electron transfer&  24 &  59 & 3.2 & $2.8 \times 10^{-8}$ \\
    Repressed & Pigment biosynthesis            &  16 &  43 & 2.9 & $6.7 \times 10^{-5}$ \\
    \bottomrule
  \end{tabular}
\end{table}
